\documentclass[11pt]{scrartcl}

\usepackage[stixtwo]{fontsetup}
\usepackage{amsmath}

\begin{document}
\begin{enumerate}
	\item $\displaystyle\lim_{x \to x_0} f(x) = A$, ha
	      $\forall \varepsilon > 0$ esetén $\exists \delta(\varepsilon) > 0$,
	      hogy $|f(x) - A| < \varepsilon$, ha $x \in D_f$ és
	      $0 < |x - x_0| < \delta(\varepsilon)$.

	\item Helyettesítsünk be a definícióban szereplő egyenlőtlenségbe:
	      \[
		      \left| \frac{10}{2x - 1} - 2 \right|
		      = \left| \frac{10 - 4x + 2}{2x - 1} \right|
		      = \frac{4 |x - 3|}{|2x - 1|}
		      \text.
	      \]
	      Becsüljük felül a törtünket. Ehhez a nevezőt kell alulról becsülni,
	      ezért szorítkozzunk a $3$ egy környezetére: ha $|x - 3| < 1$, akkor
	      $2 < x < 4$, így $|2x - 1| > 3$. Ekkor:
	      \[
		      \frac{4 |x - 3|}{|2x - 1|}
		      < \frac{4 |x - 3|}{3}
		      < \varepsilon
		      \qquad \Rightarrow \qquad
		      |x - 3| < \frac{3 \varepsilon}{4}
		      \text.
	      \]
	      A megoldás tehát:
	      \[
		      \delta(\varepsilon) = \min \left\{ 1; \frac{3 \varepsilon}{4} \right\}
		      \text.
	      \]
\end{enumerate}
\end{document}
